% sbdeck.tex -- the sample deck, the same content in every design.
% Input by B<n>.tex after \documentclass, \usepackage{sbB<n>} and the pitch
% (\sbdesignid, \sbdesignname, \sbpitchpage).
\title{Sortieren}
\subtitle{Selectionsort · Bubblesort · Mergesort}
\author{Florian Köhl}
\setsubject{Algorithmik}
\setchapter{3}
\setphotos{cover/algorithmik/playing-cards.jpg, cover/algorithmik/maze.jpg,
  cover/algorithmik/signpost.jpg, cover/algorithmik/pencils.jpg,
  cover/algorithmik/subway-map.jpg, cover/algorithmik/bridges.jpg}

\begin{document}

% 1 title slide (L1)
\sbtitleframe

% 2 pitch page
\begin{frame}{\sbdesignid\ «\sbdesignname»}
  \sbpitchpage
\end{frame}

% 3 section slide (L2)
\setcounter{section}{1}
\section{Selectionsort}
\sbsectionframe{Wiederholtes Suchen des Minimums}

% 4 bullets (L4)
\begin{frame}{Effizienz: Laufzeit}
  \begin{itemize}
    \item Wie «schnell» ist Selectionsort?
    \item Wie messen wir die Effizienz eines Algorithmus?
      \begin{itemize}
        \item Zeitmessung? Hängt von Hardware und Implementierung ab.
        \item Alternativen?
      \end{itemize}
    \item Grobe Idee: die \term{Anzahl Operationen} zählen
      \begin{itemize}
        \item Vergleiche zwischen zwei Elementen
        \item Vertauschungen zweier Elemente
      \end{itemize}
  \end{itemize}
\end{frame}

% 5 definition (L7)
\begin{frame}{Sortiert?}
  \begin{concept}[Aufsteigend sortiert]
    Eine Liste ist \term{aufsteigend sortiert}, wenn jedes Element kleiner
    oder gleich gross ist wie alle nachfolgenden Elemente.
  \end{concept}
  \begin{itemize}
    \item Sind diese Karten sortiert?
    \item Und die leere Liste?
  \end{itemize}
\end{frame}

% 6 figure (L6)
\begin{frame}{Selectionsort in Aktion}
  \vspace{5mm}
  \cardArray{name=c, items={2h,6h,5h,9h,7h,10h}, parts={1/1/S,2/6/U}, size=0.12}
  \swaparrow{c2}{c3}\pointerbelow{c3}{min}
  \vfill
  \begin{center}
    Das kleinste Element in U ist die 5:\\
    Wir vertauschen es mit dem ersten Element von U.
  \end{center}
\end{frame}

% 7 code (L10)
\begin{frame}[fragile]{Selectionsort in Python}
\begin{codebox}{selection\_sort.py}
def selection_sort(liste):
    for i in range(len(liste) - 1):
        k = i  # Position des Minimums in U
        for j in range(i + 1, len(liste)):
            if liste[j] < liste[k]:
                k = j
        liste[i], liste[k] = liste[k], liste[i]
\end{codebox}
\end{frame}

% 8 task in «Mach mit» mode (L13)
\begin{frame}[machmit={zu zweit · 15 Min}]{Ihr seid dran!}
  \vspace{2mm}
  \begin{task}{38–41}{Sortieren von Hand}
    Sortiere die Listen im Skript (S.~48) mit Selectionsort und Bubblesort
    von Hand.
  \end{task}
\end{frame}

% 9 algorithm + figure (L8)
\section{Mergesort}
\begin{frame}{Mergesort: Merge}
  \begin{columns}[T,onlytextwidth]
    \begin{column}{0.5\textwidth}
      \sbstep{1}{Vergleiche die vordersten Elemente beider Listen.}
      \sbstep[hl]{2}{Füge das kleinere hinten an die neue Liste an.}
      \sbstep{3}{Ist eine Liste leer, hänge den Rest der anderen an.}
    \end{column}
    \begin{column}{0.44\textwidth}
      \vspace{1mm}
      \pythonArray{items={2,5,3,9}, parts={0/2/links,2/4/rechts}, done={0},
        highlight={1,2}, ptrs={1/i,2/j}}
      \vspace{5mm}
      \pythonArray{items={2,\phantom{0},\phantom{0},\phantom{0}}, blank, done={0},
        drawindexes=false}
    \end{column}
  \end{columns}
\end{frame}

% 10 step by step (L9)
\section{Spannbäume}
\begin{frame}{Kruskal in Aktion}
  \begin{center}
  \begin{graph}[x=1.45cm, y=0.94cm]
    \node[vertex] (C) at (0,2) {C};
    \node[vertex] (E) at (2.6,2.5) {E};
    \node[vertex] (D) at (5,1.3) {D};
    \node[vertex] (A) at (1.3,0) {A};
    \node[vertex] (B) at (4,-0.4) {B};
    \draw (E) edge["10"] (D);
    \draw (A) edge["6"] (D);
    \draw (D) edge["5"] (B);
    \draw (A) edge["5"] (B);
    \draw (A) edge["4", onslide=<3->{removed}] (E);
    \draw (C) edge["3", onslide=<2->{chosen}] (E);
    \draw (C) edge["2", chosen] (A);
  \end{graph}
  \end{center}
  \vfill
  \begin{center}
    \only<1>{C–A (2) bildet keinen Kreis: \term{gewählt}.}%
    \only<2>{C–E (3) bildet keinen Kreis: \term{gewählt}.}%
    \only<3>{A–E (4) würde einen Kreis bilden: \term{verworfen}.}\par
    \vspace{2mm}\stepdots{3}
  \end{center}
\end{frame}

% 11 «Zum Lesen» mode
\begin{frame}[lesen={Skript S. 92}]{Ein Baum ist …}
  \begin{enumerate}
    \item ein \term{zusammenhängender} Graph ohne Kreise.
    \item ein Graph, in dem es zwischen zwei Knoten genau \term{einen Pfad} gibt.
    \item ein zusammenhängender Graph mit $n$ Knoten und $n-1$ Kanten.
  \end{enumerate}
  \medskip
  Alle drei Beschreibungen sind gleichwertig.
\end{frame}

\end{document}
